% ---- pg-173 （书页 161） ----

\chapter{多元正态分布的统计推断}
\phantomsection\label{chap:5}

\section{引言}
\phantomsection\label{sec:5.1}

本书从这一章起开始转入多元统计学的方法论. 我们将利用前面第 1 章至第 4 章的各种概念和结论建立数据分析的方法. 任何一种分析, 其大部分内容均与统计推断相联系, 即根据来自样本的信息得出有关总体的有效结论.

基于这一观点, 我们将集中讨论关于总体均值向量及其分量的统计推断问题. 虽然我们将从假设检验开始统计推断的讨论, 但我们的最终目的, 还是要基于联合置信域的形式给出均值向量诸分量的一个完整的统计分析.

多元分析的精髓之一, 就是必须对 $p$ 个相关变量同时进行分析. 我们将通过用本章中给出的方法来举例说明这一原理.

\section{\texorpdfstring{$\boldsymbol{\mu}_0$}{mu0} 作为正态总体均值的似真性}
\phantomsection\label{sec:5.2}

我们先回忆一下在一元统计理论中, 判断 $\mu_0$ 是否为总体均值 $\mu$ 的似真值的假设检验问题. 从假设检验的观点出发, 这个问题可归结为检验假设
\[
H_0: \mu = \mu_0, \qquad H_1: \mu \neq \mu_0
\]
这里 $H_0$ 为零假设, 而 $H_1$ 为(双边)备择假设. 设 $X_1, X_2, \cdots, X_n$ 为一从正态总体中抽取的随机样本, 那么, 合理的检验统计量为
\[
t = \frac{(\bar{X} - \mu_0)}{s/\sqrt{n}}
\]
其中 $\bar{X} = \dfrac{1}{n}\sum\limits_{i=1}^{n}X_i$, $s^2 = \dfrac{1}{n-1}\sum\limits_{i=1}^{n}(X_i - \bar{X})^2$.

这个检验统计量是自由度(d.f.)为 $n-1$ 的 $t$ 分布. 如果观测到的 $|t|$ 值超过自由度为 $n-1$ 的 $t$ 分布的某个特定百分位点的值, 则拒绝 $H_0$, 即 $\mu_0$ 不是 $\mu$ 的似真值.

当 $|t|$ 值较大时拒绝 $H_0$ 等价于 $t$ 的平方
\phantomsection\label{eq:5-1}
\begin{equation}
t^2 = \frac{(\bar{X} - \mu_0)^2}{s^2/n} = n(\bar{X} - \mu_0)(s^2)^{-1}(\bar{X} - \mu_0) \tag{5-1}
\end{equation}
较大时拒绝 $H_0$. 在\hyperref[eq:5-1]{式(5-1)}中, 变量 $t^2$ 为样本均值 $\bar{X}$ 到检验值 $\mu_0$ 的距离的平方. 距离的单位可用 $s/\sqrt{n}$ 或者 $\bar{X}$ 的估计标准差来表示. 一旦观测出 $\bar{X}$ 和 $s^2$, 则检验的原则为: 在显著性水平 $\alpha$ 下, 如果
\phantomsection\label{eq:5-2}
\begin{equation}
n(\bar{x} - \mu_0)(s^2)^{-1}(\bar{x} - \mu_0) > t_{n-1}^2(\alpha/2) \tag{5-2}
\end{equation}
则拒绝 $H_0$ 而接受 $H_1$, 其中 $t_{n-1}(\alpha/2)$ 代表自由度为 $n-1$ 的 $t$ 分布的上 $100(\alpha/2)$ 百分位数.

如果 $H_0$ 没有被拒绝, 那么我们就认为 $\mu_0$ 为正态总体均值的似真值. 然而, 问题是, 是否还
